\uuid{9qB5}
\titre{Loi binomiale et approximation normale }

\niveau{}
\module{ Probabilités }
\chapitre{Loi binomiale}
\sousChapitre{Loi binomiale et approximation}
\theme{Probabilité et loi binomiale}
\auteur{Quentin Liard}
\datecreate{ 2026-04-03}
\organisation{AMSCC}
\difficulte{3}

\contenu{

\texte{ 
Une entreprise fabrique en grande série des capteurs. La probabilité qu’un capteur soit défectueux est $p=0{,}08$.
On prélève au hasard un lot de $n=200$ capteurs, en supposant les fabrications indépendantes.
On note $X$ le nombre de capteurs défectueux dans le lot.

}

\begin{enumerate}

\item   \question{Justifier que $X$ suit une loi binomiale, dont on précisera les paramètres.}
\indication{}
\reponse{
Chaque capteur correspond à une épreuve de Bernoulli indépendante, avec succès = ``le capteur est défectueux'' de probabilité $p=0{,}08$. Le nombre de succès sur $n=200$ épreuves suit donc une loi binomiale :
\[
X\sim \mathcal{B}(200 ; 0{,}08).
\]
}

\item   \question{Calculer l’espérance et la variance de $X$.}
\indication{}
\reponse{
Pour une loi binomiale $\mathcal{B}(n,p)$,
\[
E(X)=np,\qquad V(X)=np(1-p).
\]
Donc ici
\[
E(X)=200\times 0{,}08=16,
\]
et
\[
V(X)=200\times 0{,}08\times 0{,}92=14{,}72.
\]
}

\item   \question{Déterminer une valeur approchée de la probabilité $P(X=20)$.}
\indication{}
\reponse{
\[
P(X=20)=\binom{200}{20}(0{,}08)^{20}(0{,}92)^{180}.
\]
Numériquement,
\[
P(X=20)\approx 0{,}0564.
\]
}

%\item   \question{Calculer la probabilité $P(X\geq 15)$.}
%\indication{}
%\reponse{
%\[
%P(X\geq 15)=\sum_{k=15}^{200}\binom{200}{k}(0{,}08)^k(0{,}92)^{200-k}.
%\]
%Numériquement,
%\[
%P(X\geq 15)\approx 0{,}6403.
%\]
%}

\item   \question{En approchant $X$ par une loi normale $Y$ de même espérance et de même variance, estimer $P(Y\leq 10)$.}
\indication{}
\reponse{
On approche $X$ par une variable aléatoire
\[
Y\sim \mathcal{N}(16\,;\,14{,}72).
\]
Alors
\[
P(X\leq 10)\approx P(Y\leq 10).
\]
On standardise :
\[
P(Y\leq 10)=
P\left(\frac{Y-16}{\sqrt{14{,}72}}\leq \frac{10-16}{\sqrt{14{,}72}}\right).
\]
Or
\[
\frac{10-16}{\sqrt{14{,}72}}\approx -1{,}56.
\]
Donc
\[
P(X\leq 10)\approx \Phi(-1{,}56)\approx 0{,}059.
\]
}

%\item   \question{En approchant $X$ par une loi normale de même espérance et de même variance, estimer $P(12\leq X\leq 20)$.}
%\indication{}
%\reponse{
%On approche $X$ par une variable aléatoire
%\[
%Y\sim \mathcal{N}(16\,;\,14{,}72).
%\]
%Alors
%\[
%P(12\leq X\leq 20)\approx P(12\leq Y\leq 20).
%\]
%On standardise :
%\[
%P(12\leq Y\leq 20)
%=
%P\left(\frac{12-16}{\sqrt{14{,}72}}\leq \frac{Y-16}{\sqrt{14{,}72}}%\leq \frac{20-16}{\sqrt{14{,}72}}\right).
%\]
%Or
%\[
%\frac{12-16}{\sqrt{14{,}72}}\approx -1{,}04,
%\qquad
%\frac{20-16}{\sqrt{14{,}72}}\approx 1{,}04.
%\]
%Donc
%\[
%P(12\leq X\leq 20)\approx \Phi(1{,}04)-\Phi(-1{,}04)\approx 0{,}702.
%\]
%}

\end{enumerate}
}